% Formation Engine -- XYZ File Formats (Edition 2)
% VSEPR-SIM 4.0.0  |  Compile: pdflatex (twice)
%
\documentclass[12pt, a4paper, oneside]{article}

\usepackage[top=1.0in, bottom=1.0in, left=1.15in, right=1.15in, headheight=14pt]{geometry}
\usepackage[utf8]{inputenc}
\usepackage[T1]{fontenc}
\usepackage{amsmath, amssymb, mathtools}
\usepackage{booktabs}
\usepackage{array}
\usepackage{tabularx}
\usepackage{enumitem}
\setlist{nosep, leftmargin=1.5em}
\usepackage{listings}
\usepackage{xcolor}

\lstset{basicstyle=\ttfamily\small, keywordstyle=\color{blue}, commentstyle=\color{gray},
  stringstyle=\color{red}, breaklines=true, columns=fullflexible, keepspaces=true,
  frame=single, backgroundcolor=\color{gray!8}, xleftmargin=0.5em, xrightmargin=0.5em,
  language=C++}

\usepackage{graphicx}
\usepackage{fancyhdr}
\pagestyle{fancy}
\fancyhf{}
\fancyhead[L]{\small\textit{Formation Engine Methodology}}
\fancyhead[R]{\small\textit{XYZ File Formats --- Edition 2}}
\fancyfoot[C]{\thepage}
\renewcommand{\headrulewidth}{0.4pt}

\usepackage{titlesec}
\titleformat{\section}{\Large\bfseries}{\S{}\thesection}{1em}{}
\titleformat{\subsection}{\large\bfseries}{\thesubsection}{1em}{}
\titleformat{\subsubsection}{\normalsize\bfseries}{\thesubsubsection}{1em}{}

\usepackage[colorlinks=true, linkcolor=black, citecolor=black, urlcolor=blue]{hyperref}

\newcommand{\file}[1]{\texttt{#1}}
\newcommand{\type}[1]{\texttt{#1}}
\newcommand{\field}[1]{\texttt{#1}}
\newcommand{\Ang}{\text{\AA}}
\newcommand{\kcalmol}{\text{kcal/mol}}

\title{%
  \textbf{Formation Engine Methodology}\\[6pt]
  \textbf{XYZ File Formats}\\[10pt]
  \large VSEPR-SIM Atomistic Engine\\
  Formats: \texttt{.xyz}\quad\texttt{.xyza}\quad\texttt{.xyzc}\quad\texttt{.xyzf}\\[4pt]
  \normalsize Edition 2 --- Expanded Reference
}
\author{VSEPR-SIM Project}
\date{Version 4.0.0 --- \today}

\begin{document}
\maketitle

\begin{abstract}
All VSEPR-SIM molecular I/O uses plain-text, line-oriented files.
Four formats share a common structural grammar; each adds fields to
the preceding one.  This document is the expanded second-edition
reference for the \texttt{.xyz}, \texttt{.xyza}, \texttt{.xyzc}, and
\texttt{.xyzf} file formats, including grammar specifications, field
semantics, coordinate conventions, checkpoint and trajectory usage,
bounding-box encoding, implementation mapping, and typical observable
ranges.
\end{abstract}

\tableofcontents
\newpage

% ================================================
\section{Overview}
% ================================================

The VSEPR-SIM atomistic engine reads and writes molecular geometry,
dynamics, and checkpoint data using four progressively richer
plain-text formats.  Each format extends the previous one by appending
per-atom or per-frame metadata while preserving backward compatibility
with simple line-oriented parsers.

\begin{table}[htbp]
  \centering
  \caption{Format hierarchy.}
  \label{tab:hierarchy}
  \begin{tabular}{llp{8cm}}
    \toprule
    Extension & Name & Use case \\
    \midrule
    \texttt{.xyz}  & Standard XYZ     & Static geometry snapshot \\
    \texttt{.xyza} & Extended XYZ     & Charges, velocities, forces \\
    \texttt{.xyzc} & Checkpoint XYZ   & Restartable MD state \\
    \texttt{.xyzf} & Formation XYZ    & Multi-frame trajectory \\
    \bottomrule
  \end{tabular}
\end{table}

\subsection{Unit Conventions}

All formats share a single, non-negotiable unit system:

\begin{table}[htbp]
  \centering
  \caption{Unit system (all formats).}
  \label{tab:units}
  \begin{tabular}{lll}
    \toprule
    Quantity & Unit & SI equivalent \\
    \midrule
    Coordinates   & \AA{} (angstrom) & $10^{-10}$\,m \\
    Energy        & kcal/mol         & $4.184 \times 10^{3}$\,J/mol \\
    Velocity      & \AA/fs           & $10^{5}$\,m/s \\
    Force         & kcal/(mol\,\AA)  & $6.947 \times 10^{13}$\,N/mol \\
    Charge        & elementary $e$   & $1.602 \times 10^{-19}$\,C \\
    Time          & fs               & $10^{-15}$\,s \\
    Temperature   & K                & kelvin \\
    \bottomrule
  \end{tabular}
\end{table}

\subsection{Coordinate System}

Right-handed Cartesian $(x, y, z)$ in \AA.  Origin is arbitrary;
the conventional choice is near the centre of mass.  Coordinate
precision: six decimal places ($\Rightarrow$ resolution $\approx 0.1$\,pm).

% ================================================
\section{Standard XYZ (\texttt{.xyz})}
\label{sec:xyz}
% ================================================

\subsection{Grammar}

\begin{lstlisting}[language={},frame=single]
<N>
<comment>
<sym> <x> <y> <z>
...  (N lines)
\end{lstlisting}

\begin{itemize}
  \item \textbf{Line 1}: atom count $N$ (integer $\ge 1$).
  \item \textbf{Line 2}: free-text comment (may be empty).
  \item \textbf{Lines 3--$(N{+}2)$}: element symbol + three Cartesian
        coordinates, whitespace-separated.
\end{itemize}

\subsection{Comment Line Convention}

The comment line follows a structured-but-optional pattern:

\begin{lstlisting}[language={},frame=single]
<name> | E = <val> kcal/mol | <notes>
\end{lstlisting}

Fields are pipe-delimited.  The energy key \texttt{E =} is recognised
by the reader for energy-tracking contexts.  Additional free-text
notes are permitted after the second pipe.

\subsection{Example: Water}

\begin{lstlisting}[language={},frame=single]
3
H2O | E = -6.245 kcal/mol
O   0.000000   0.000000   0.119262
H   0.000000   0.763239  -0.477049
H   0.000000  -0.763239  -0.477049
\end{lstlisting}

\subsection{Parsing Notes}

\begin{itemize}
  \item Element symbols are case-sensitive: \texttt{He} is helium,
        \texttt{HE} is rejected.
  \item Trailing whitespace on atom lines is ignored.
  \item Blank lines between the comment line and the first atom line
        are \emph{not} permitted; the reader expects $N$ consecutive
        atom lines starting at line 3.
  \item Coordinate fields must parse as IEEE~754 double-precision
        floating-point values.
\end{itemize}

% ================================================
\section{Extended XYZ (\texttt{.xyza})}
\label{sec:xyza}
% ================================================

Adds per-atom scalar and vector properties after the $(x, y, z)$ fields.

\subsection{Grammar}

\begin{lstlisting}[language={},frame=single]
<N>
<comment> properties="<list>","<sym> <x> <y> <z> [<q> <vx> <vy> <vz> <fx> <fy> <fz> <e>]
\end{lstlisting}

The \texttt{properties} key declares which optional columns are present.
Column order is fixed; omission is permitted only from the right.

\subsection{Property Columns}

\begin{table}[htbp]
  \centering
  \caption{Extended property columns (fixed order).}
  \label{tab:xyza_cols}
  \begin{tabular}{llll}
    \toprule
    Key & Symbol & Unit & Columns \\
    \midrule
    \texttt{charge}   & $q_i$          & $e$                   & 1 \\
    \texttt{velocity} & $\mathbf{v}_i$ & \AA/fs                & 3 \\
    \texttt{force}    & $\mathbf{F}_i$ & kcal/(mol\,\AA)       & 3 \\
    \texttt{energy}   & $\varepsilon_i$& kcal/mol              & 1 \\
    \bottomrule
  \end{tabular}
\end{table}

\subsection{Properties Declaration Syntax}

The \texttt{properties} value is a colon-separated list of keys:

\begin{lstlisting}[language={},frame=single]
properties="charge:velocity"
properties="charge:velocity:force:energy"
\end{lstlisting}

Unrecognised keys trigger a parse warning but do not abort the read.

\subsection{Example: H\texorpdfstring{$_2$}{2}O with Charges and Velocities}

\begin{lstlisting}[language={},frame=single]
3
Water MD frame 1 properties="charge:velocity"
O   0.00   0.00   0.12  -0.834   0.00   0.00   0.00
H   0.00   0.76  -0.48   0.417   0.12   0.08  -0.03
H   0.00  -0.76  -0.48   0.417  -0.12   0.08  -0.03
\end{lstlisting}

\subsection{Charge Semantics}

The charge field $q_i$ represents the effective partial charge used in
Coulomb interactions.  It is \emph{not} an oxidation state.  For
Drude-polarisable models, $q_i$ is the core charge; the Drude
particle charge is stored separately in the polarisation subsystem.

% ================================================
\section{Checkpoint XYZ (\texttt{.xyzc})}
\label{sec:xyzc}
% ================================================

Extends \texttt{.xyza} with a simulation-state header block required
to restart an interrupted MD run without loss of trajectory continuity.

\subsection{Header Block}

The header block appears \emph{before} the atom-count line, delimited
by \texttt{CHECKPOINT} / \texttt{END\_CHECKPOINT} markers:

\begin{lstlisting}[language={},frame=single]
CHECKPOINT
step       <integer>
time       <float> fs
dt         <float> fs
T_target   <float> K
seed       <integer>
box        <ax> <ay> <az>
END_CHECKPOINT
<N>
<comment>
<atom lines with charge + velocity + force>
\end{lstlisting}

\subsection{Header Field Semantics}

\begin{table}[htbp]
  \centering
  \caption{Checkpoint header fields.}
  \label{tab:xyzc_header}
  \begin{tabular}{llp{7.5cm}}
    \toprule
    Field & Type & Meaning \\
    \midrule
    \texttt{step}      & integer & MD step number at checkpoint \\
    \texttt{time}      & float   & Accumulated simulation time (fs) \\
    \texttt{dt}        & float   & Integrator timestep (fs) \\
    \texttt{T\_target} & float   & Thermostat set-point (K); 0 = NVE ensemble \\
    \texttt{seed}      & integer & RNG seed for reproducible continuation \\
    \texttt{box}       & 3 floats & Simulation box side lengths (\AA);
                                    omit line entirely for open boundary \\
    \bottomrule
  \end{tabular}
\end{table}

\subsection{Restart Protocol}

When the engine loads a \texttt{.xyzc} file:

\begin{enumerate}
  \item The header is parsed; integrator state restored (step, time, dt,
        thermostat target, RNG state).
  \item The atom block is read as a standard \texttt{.xyza} frame.
  \item Neighbour lists and bond tables are reconstructed from coordinates;
        they are \emph{never} serialised into the checkpoint.
  \item If a \texttt{box} line is present, PBC is activated.
\end{enumerate}

\subsection{Checkpoint vs.\ Trajectory}

A \texttt{.xyzc} file contains exactly one frame.  For multi-frame
output, use \texttt{.xyzf} (\S\ref{sec:xyzf}).

% ================================================
\section{Formation XYZ (\texttt{.xyzf})}
\label{sec:xyzf}
% ================================================

A multi-frame format: one or more \texttt{.xyz} or \texttt{.xyza} blocks
concatenated in a single file.  Used for minimisation traces, NVT
trajectories, and reaction paths.

\subsection{Grammar}

\begin{lstlisting}[language={},frame=single]
<N>           <- frame 1
<comment>
<atom lines>
<N>           <- frame 2
<comment>
<atom lines>
...
\end{lstlisting}

Each frame is self-contained and independently parseable.

\subsection{Comment Line Convention (Multi-Frame)}

The comment line of frame $k$ typically records the step index and
instantaneous energy:

\begin{lstlisting}[language={},frame=single]
Step 0 | E = -12.30 kcal/mol | FIRE iter
Step 1 | E = -13.11 kcal/mol | FIRE iter
\end{lstlisting}

\subsection{Frame Independence}

Each frame must be independently parseable:

\begin{itemize}
  \item Atom count $N$ may vary between frames.
  \item The \texttt{properties} declaration may differ between frames.
  \item Readers must re-parse the header of each frame; caching the
        first frame's column layout globally is a bug.
\end{itemize}

\subsection{Typical Uses}

\begin{table}[htbp]
  \centering
  \caption{Common \texttt{.xyzf} applications.}
  \label{tab:xyzf_uses}
  \begin{tabular}{lp{8cm}}
    \toprule
    Application & Comment line pattern \\
    \midrule
    FIRE minimisation trace  & \texttt{Step <k> | E = <val> kcal/mol | FIRE iter} \\
    NVT trajectory dump      & \texttt{Step <k> | E = <val> kcal/mol | T = <T> K} \\
    Reaction path (NEB)      & \texttt{Image <k> | E = <val> kcal/mol | dE = <dE>} \\
    Conformer scan           & \texttt{Conf <k> | E = <val> kcal/mol | <dihedral>} \\
    \bottomrule
  \end{tabular}
\end{table}

% ================================================
\section{Bounding Box and Periodic Boundary Conditions}
\label{sec:pbc}
% ================================================

When PBC are active, box geometry is stored in the \texttt{.xyzc} header
(\S\ref{sec:xyzc}) or as comment-line keys in any format.

\subsection{Comment-Line Encoding}

\begin{lstlisting}[language={},frame=single]
Lattice="ax 0 0  0 ay 0  0 0 az" pbc="T T T"
\end{lstlisting}

\begin{itemize}
  \item Box vectors in \AA, flattened $3\times 3$ row-major matrix.
        Orthorhombic: diagonal entries only; triclinic: full matrix.
  \item \texttt{pbc} flags: \texttt{T} = periodic, \texttt{F} = open.
  \item Compatible with ASE Extended XYZ convention.
\end{itemize}

\subsection{Box Vector Interpretation}

For an orthorhombic cell with side lengths $a_x, a_y, a_z$:
\begin{equation}
  \mathbf{L} =
  \begin{pmatrix}
    a_x & 0   & 0   \\
    0   & a_y & 0   \\
    0   & 0   & a_z
  \end{pmatrix}
  \label{eq:lattice}
\end{equation}
The cell volume is $V = \det(\mathbf{L}) = a_x \cdot a_y \cdot a_z$.

\subsection{Minimum Image Convention}

Under PBC, the engine applies the minimum-image convention:
\begin{equation}
  \Delta r_\alpha = r_\alpha^{(j)} - r_\alpha^{(i)}
    - L_\alpha \cdot \mathrm{round}\!\left(
      \frac{r_\alpha^{(j)} - r_\alpha^{(i)}}{L_\alpha}
    \right),
  \qquad \alpha \in \{x, y, z\}
  \label{eq:mic}
\end{equation}
where $L_\alpha$ is the box length along axis $\alpha$.

% ================================================
\section{Bond Auto-Detection}
\label{sec:bonds}
% ================================================

Bond connectivity is inferred from geometry, never stored in the file.
The \texttt{detect\_bonds} function uses per-element covalent radii:
\begin{equation}
  d_{ij} < s \cdot \bigl(r_{\text{cov}}(Z_i) + r_{\text{cov}}(Z_j)\bigr)
  \quad \Longrightarrow \quad \text{bond}(i, j)
  \label{eq:bond}
\end{equation}
where $s$ is the scale factor (default $s = 1.2$).

\subsection{Design Rationale}

Bonds are derived data, not ground truth.  Stale bond tables are a
silent-corruption vector.  The \S{}0 container philosophy formalises
this: caches never become truth.

% ================================================
\section{Implementation Mapping}
\label{sec:impl}
% ================================================

\begin{table}[htbp]
  \centering
  \caption{Source files implementing the XYZ format family.}
  \label{tab:impl}
  \begin{tabular}{lp{8.5cm}}
    \toprule
    File & Role \\
    \midrule
    \file{include/io/xyz\_format.hpp}
      & Public API: \type{XYZReader}, \type{XYZWriter},
        \type{XYZAReader}, \type{XYZAWriter} \\
    \file{src/io/xyz\_format.cpp}
      & Reader / writer implementation \\
    \file{include/thermal/xyzc\_format.hpp}
      & Checkpoint reader/writer declarations \\
    \file{src/thermal/xyzc\_format.cpp}
      & Checkpoint reader/writer implementation \\
    \bottomrule
  \end{tabular}
\end{table}

All reader and writer types reside in namespace \type{vsepr::io}.

\subsection{Reader API Pattern}

\begin{lstlisting}
// Standard XYZ
auto frame = vsepr::io::XYZReader::read("molecule.xyz");
// Extended XYZ
auto frame = vsepr::io::XYZAReader::read("dynamics.xyza");
// Checkpoint
auto state = vsepr::io::XYZCReader::read("restart.xyzc");
// Formation (multi-frame)
auto traj  = vsepr::io::XYZFReader::read_all("trajectory.xyzf");
// or frame-by-frame:
vsepr::io::XYZFReader reader("trajectory.xyzf");
while (auto frame = reader.next()) { /* ... */ }
\end{lstlisting}

\subsection{Writer API Pattern}

\begin{lstlisting}
vsepr::io::XYZWriter::write("output.xyz", atoms, comment);
vsepr::io::XYZAWriter::write("output.xyza", atoms, props, comment);
vsepr::io::XYZCWriter::write("checkpoint.xyzc", state);
vsepr::io::XYZFWriter writer("trajectory.xyzf", /*append=*/true);
writer.write_frame(atoms, comment);
\end{lstlisting}

% ================================================
\section{Typical Energy and Observable Ranges}
\label{sec:ranges}
% ================================================

\begin{table}[htbp]
  \centering
  \caption{Typical observable magnitudes for sanity checking.}
  \label{tab:ranges}
  \begin{tabular}{llr}
    \toprule
    Observable & Typical range & Unit \\
    \midrule
    Total bonded energy          & $-10^2$ to $-10^5$   & kcal/mol \\
    Relative conformer energy    & 0.1 -- 10            & kcal/mol \\
    Activation barrier           & 5 -- 50              & kcal/mol \\
    Force (unrelaxed atom)       & 1 -- 100             & kcal/(mol\,\AA) \\
    Force convergence criterion  & $< 0.01$             & kcal/(mol\,\AA) \\
    Velocity (300\,K, C atom)    & $\sim 0.02$          & \AA/fs \\
    Partial charge (O in water)  & $-0.82$ to $-0.84$   & $e$ \\
    Partial charge (H in water)  & $+0.41$ to $+0.42$   & $e$ \\
    \bottomrule
  \end{tabular}
\end{table}

These ranges serve as first-pass sanity checks.  Values outside these
bounds in production output should trigger audit warnings.

% ================================================
\section{Integration with the Identity--State Framework}
\label{sec:s0}
% ================================================

The XYZ formats serve as the serialisation layer for the particle
identity vector $\mathbf{I}_i$ defined in \S{}0 of the Formation
Engine methodology.

\begin{table}[htbp]
  \centering
  \caption{Identity vector components and their file-format encoding.}
  \label{tab:s0_map}
  \begin{tabular}{llp{6.5cm}}
    \toprule
    Component & Symbol & File encoding \\
    \midrule
    Nuclear identity     & $Z_i$       & Element symbol (looked up to $Z$) \\
    Mass identity        & $A_i$       & Implicit via $Z_i$ \\
    Charge participation & $Q_i$       & \texttt{charge} column in \texttt{.xyza} \\
    Structural role      & $\Sigma_i$  & Not serialised; inferred at load \\
    Stability class      & $\Lambda_i$ & Not serialised; computed post-load \\
    Provenance memory    & $\Theta_i$  & Comment-line metadata \\
    \bottomrule
  \end{tabular}
\end{table}

$\mathcal{K}_{\text{geom}}$ maps to the lattice matrix from the
$\texttt{Lattice}$ key or \texttt{box} header.
$\mathcal{K}_{\text{comp}}$ is reconstructed at load from element
composition and geometry.  The source-layer flag $F_K \in \{Z, A, C\}$
maps \texttt{Z}$\to$\texttt{.xyz}, \texttt{A}$\to$\texttt{.xyza},
\texttt{C}$\to$\texttt{.xyzc}.

% ================================================
\section{Format Comparison Matrix}
\label{sec:matrix}
% ================================================

\begin{table}[htbp]
  \centering
  \caption{Feature comparison across formats.}
  \label{tab:matrix}
  \begin{tabular}{lcccc}
    \toprule
    Feature & \texttt{.xyz} & \texttt{.xyza} & \texttt{.xyzc} & \texttt{.xyzf} \\
    \midrule
    Atom count + coordinates       & \checkmark & \checkmark & \checkmark & \checkmark \\
    Partial charges                &            & \checkmark & \checkmark & opt.       \\
    Velocities                     &            & \checkmark & \checkmark & opt.       \\
    Forces                         &            & \checkmark & \checkmark & opt.       \\
    Per-atom energy                &            & \checkmark & \checkmark & opt.       \\
    Integrator state               &            &            & \checkmark &            \\
    Thermostat / RNG               &            &            & \checkmark &            \\
    Box geometry (header)          &            &            & \checkmark &            \\
    Box geometry (comment key)     & \checkmark & \checkmark & \checkmark & \checkmark \\
    Multi-frame                    &            &            &            & \checkmark \\
    \bottomrule
  \end{tabular}
\end{table}

% ================================================
\section{Error Handling and Validation}
\label{sec:errors}
% ================================================

Readers enforce these validation rules at parse time:

\begin{enumerate}
  \item \textbf{Atom count mismatch.}  Fewer than $N$ atom lines
        $\Rightarrow$ \type{XYZParseError::INCOMPLETE\_FRAME}.
  \item \textbf{Unknown element.}  Unrecognised symbol
        $\Rightarrow$ \type{XYZParseError::UNKNOWN\_ELEMENT}.
        Two-letter symbols must match the periodic table (case-sensitive).
  \item \textbf{Malformed coordinates.}  Non-numeric field
        $\Rightarrow$ \type{XYZParseError::BAD\_COORDINATE}.
  \item \textbf{Column count mismatch.}  Declared \texttt{properties}
        implies $k$ columns but line has fewer: warn and zero-fill
        (fail-soft for velocities/forces; fail-hard for coordinates).
  \item \textbf{Checkpoint integrity.}  Missing
        \texttt{CHECKPOINT}/\texttt{END\_CHECKPOINT} block
        $\Rightarrow$ \type{XYZCParseError::NO\_HEADER}.
\end{enumerate}

% ================================================
\section{Cross-References}
% ================================================

\begin{table}[htbp]
  \centering
  \caption{Related documentation.}
  \label{tab:xrefs}
  \begin{tabular}{lp{8cm}}
    \toprule
    Document & Relevance \\
    \midrule
    \file{docs/FILE\_FORMATS.md}              & Prose description of all I/O formats \\
    \file{section0\_identity\_state.pdf}       & Formal particle/cell/world definitions \\
    \file{section5\_integration.tex}           & MD timestep and integrator details \\
    \file{ALPHA\_MODEL\_BOOKLET.tex}           & Alpha model (Drude charge coupling) \\
    \file{section\_fuzzy\_ball\_model.tex}     & Coarse-grained bead visualisation \\
    \bottomrule
  \end{tabular}
\end{table}

\vfill

\begin{center}
\rule{0.5\textwidth}{0.4pt}\\[6pt]
\small Formation Engine Methodology v3.0.0 --- XYZ File Formats (Edition 2)
\end{center}

\end{document}
